\chapter{점과 화살표}

공책에는 줄이 수십만 개나 있어요.
줄을 하나하나 읽어서는 누가 믿을 만한지 한눈에 보이지 않아요.
그래서 그림으로 바꿔 보려고 해요.

규칙은 간단해요.

\begin{enumerate}
\item 통마다 점을 하나씩 찍어요. 사람의 통은 동그라미, 로봇의 통은 둥근 네모로 그려요.
\item 공책에 “가가 나에게 보냈다”는 줄이 있으면, 가의 점에서 나의 점으로 화살표를 그려요.
\item 같은 두 점 사이에 줄이 여러 개 있으면 화살표를 여러 개 그리지 않고, 하나만 그려서 그 위에 횟수나 양을 적어요.
\end{enumerate}

이렇게 점과 화살표로 그린 그림을 수학에서는 \textbf{그래프}라고 불러요.
막대그래프나 꺾은선그래프와는 다른, 이름만 같은 그림이에요.
점은 \textbf{노드}, 화살표는 \textbf{엣지}라고도 해요.

\section{한 교실, 두 장의 지도}

우리 반의 기록을 그려 볼게요.
송금 줄로 그리면 왼쪽 지도가 되고, 허락 줄로 그리면 오른쪽 지도가 돼요.

\begin{figure}[h]
\centering
\begin{tikzpicture}[scale=0.72, every node/.style={transform shape}]
  \begin{scope}
    \node[kid] (a) at (90:2.3) {지우};
    \node[kid] (b) at (162:2.3) {하준};
    \node[kid] (c) at (234:2.3) {서연};
    \node[kid] (d) at (306:2.3) {도윤};
    \node[kid] (e) at (18:2.3) {민서};
    \node[robot] (r) at (0,0) {로봇};
    \draw[send] (a) -- (b);
    \draw[send] (b) to[bend right=12] (e);
    \draw[send] (d) -- (c);
    \draw[send] (c) to[bend left=25] (a);
    \draw[send] (e) -- (d);
    \draw[send] (a) -- (r);
    \draw[send] (b) -- (r);
    \draw[send] (c) -- (r);
    \draw[send] (r) -- (d);
    \draw[send] (e) to[bend right=12] (b);
    \draw[send] (d) to[bend left=10] (b);
    \node[font=\bfseries, color=sendcol] at (0,-3.4) {송금 지도};
  \end{scope}
  \begin{scope}[xshift=7.2cm]
    \node[kid] (a) at (90:2.3) {지우};
    \node[kid] (b) at (162:2.3) {하준};
    \node[kid] (c) at (234:2.3) {서연};
    \node[kid] (d) at (306:2.3) {도윤};
    \node[kid] (e) at (18:2.3) {민서};
    \node[robot] (r) at (0,0) {로봇};
    \draw[allow] (a) -- (r);
    \draw[allow] (b) -- (r);
    \draw[allow] (c) -- (r);
    \draw[allow] (e) -- (r);
    \draw[allow] (d) -- (c);
    \node[font=\bfseries, color=allowcol] at (0,-3.4) {허락 지도};
  \end{scope}
\end{tikzpicture}
\caption{같은 교실을 송금 줄로 그린 지도(왼쪽)와 허락 줄로 그린 지도(오른쪽).}
\end{figure}

두 지도를 나란히 놓고 보면 금방 차이가 보여요.
송금 지도에는 화살표가 이리저리 얽혀 있어요.
허락 지도는 훨씬 단순해요.
화살표가 적고, 대부분 가운데 로봇을 향해 있어요.
3장에서 말한 그대로예요.
허락은 드물고, 주로 로봇이 받아요.

\section{가장 쉬운 점수: 받은 화살표 세기}

이제 첫 번째 평판 점수를 만들어 볼게요.
가장 쉬운 방법은 \emph{들어오는 화살표의 수}를 세는 거예요.
많은 친구에게서 허락을 받은 통은 많은 친구가 믿는 통이니까요.

\begin{center}
\small
\begin{tabular}{@{}lcc@{}}
\toprule
& 송금 지도에서 받은 화살표 & 허락 지도에서 받은 화살표 \\
\midrule
지우 & 1 & 0 \\
하준 & 3 & 0 \\
서연 & 1 & 1 \\
도윤 & 2 & 0 \\
민서 & 1 & 0 \\
로봇 & 3 & 4 \\
\bottomrule
\end{tabular}
\end{center}

송금 지도에서는 하준이와 로봇이 3개로 공동 1등이에요.
허락 지도에서는 로봇이 4개로 1등이고, 서연이가 1개로 2등이에요.
같은 교실인데 어느 지도를 보느냐에 따라 1등이 달라졌어요.
\emph{무엇을} 세느냐가 점수의 뜻을 바꾸는 거예요.

받은 화살표 수를 세는 방법은 단순하지만 생각보다 쓸모가 많아요.
13장에서 이 단순한 방법이 아주 강한 맞수로 다시 나타나요.
그런데 이 방법에는 큰 구멍이 하나 있어요.
그 구멍은 다음 장에서 이야기할게요.

\begin{deeper}[그래프의 역사]
점과 선으로 문제를 푸는 생각은 1736년 수학자 레온하르트 오일러에게서 시작되었다고 해요.
쾨니히스베르크라는 도시에 다리가 일곱 개 있었는데,
“모든 다리를 한 번씩만 건너서 산책할 수 있을까?”라는 수수께끼가 있었어요.
오일러는 땅을 점으로, 다리를 선으로 바꾸어 그 산책이 불가능하다는 것을 보였어요.
이 책의 지도는 화살표에 방향이 있으므로 \textbf{방향 그래프}라고 불러요.
들어오는 화살표의 수는 \textbf{진입 차수}라고 해요.
논문에서는 허락 지도의 진입 차수를 “받은 허락 차수”, 송금 지도의 진입 차수를 “받은 송금 차수”라는 뜻의 영어 이름으로 불렀어요.
\end{deeper}

\begin{learned}
\begin{itemize}
\item 공책의 줄은 점과 화살표로 된 지도(그래프)로 바꿀 수 있어요.
\item 같은 교실도 송금 지도와 허락 지도가 아주 달라요.
\item 가장 쉬운 점수는 받은 화살표의 수를 세는 거예요.
\end{itemize}
\end{learned}
